\nonstopmode{}
\documentclass[letterpaper]{book}
\usepackage[times,inconsolata,hyper]{Rd}
\usepackage{makeidx}
\makeatletter\@ifl@t@r\fmtversion{2018/04/01}{}{\usepackage[utf8]{inputenc}}\makeatother
% \usepackage{graphicx} % @USE GRAPHICX@
\makeindex{}
\begin{document}
\chapter*{}
\begin{center}
{\textbf{\huge Package `geomattR'}}
\par\bigskip{\large \today}
\end{center}
\ifthenelse{\boolean{Rd@use@hyper}}{\hypersetup{pdftitle = {geomattR: Calculate geometric attributes of spatial polygons}}}{}
\ifthenelse{\boolean{Rd@use@hyper}}{\hypersetup{pdfauthor = {Gabriel Ortega-Solis}}}{}
\begin{description}
\raggedright{}
\item[Type]\AsIs{Package}
\item[Title]\AsIs{Calculate geometric attributes of spatial polygons}
\item[Version]\AsIs{0.1.0}
\item[Description]\AsIs{What the package does (one paragraph).}
\item[License]\AsIs{MIT + file LICENSE}
\item[Encoding]\AsIs{UTF-8}
\item[Roxygen]\AsIs{list(markdown = TRUE)}
\item[RoxygenNote]\AsIs{7.3.3}
\item[Imports]\AsIs{terra, geosphere, methods}
\item[Suggests]\AsIs{testthat (>= 3.0.0)}
\item[Config/testthat/edition]\AsIs{3}
\item[NeedsCompilation]\AsIs{no}
\item[Author]\AsIs{Gabriel Ortega-Solis [aut, cre] (ORCID:
<}\url{https://orcid.org/0000-0002-0516-5694}\AsIs{>)}
\item[Maintainer]\AsIs{Gabriel Ortega-Solis }\email{g.ortega.solis@gmail.com}\AsIs{}
\end{description}
\Rdcontents{Contents}
\HeaderA{geomattR-package}{geomattR: Calculate geometric attributes of spatial polygons}{geomattR.Rdash.package}
\aliasA{geomattR}{geomattR-package}{geomattR}
\keyword{internal}{geomattR-package}
%
\begin{Description}
What the package does (one paragraph).
\end{Description}
%
\begin{Author}
\strong{Maintainer}: Gabriel Ortega-Solis \email{g.ortega.solis@gmail.com} (\Rhref{https://orcid.org/0000-0002-0516-5694}{ORCID})

\end{Author}
\HeaderA{calculate\_geometric\_attributes}{Calculate Geometric Attributes of Spatial Polygons}{calculate.Rul.geometric.Rul.attributes}
%
\begin{Description}
Calculates a comprehensive set of geometric and morphometric attributes for
spatial polygon features including area, perimeter, compactness metrics,
shape indices, elongation, and orientation measures. If multiple features are
provided, processes each feature sequentially.
\end{Description}
%
\begin{Usage}
\begin{verbatim}
calculate_geometric_attributes(v, metrics = "all")
\end{verbatim}
\end{Usage}
%
\begin{Arguments}
\begin{ldescription}
\item[\code{v}] A SpatVector object representing one or more polygons.

\item[\code{metrics}] Character string or vector specifying which metrics to calculate.
Options are:
\begin{itemize}

\item{} \code{"all"} (default): Calculate all available metrics
\item{} A single metric name as a string: Calculate only that metric
\item{} A character vector of metric names: Calculate specified metrics

\end{itemize}

Available metric names: "area", "perimeter", "compactness", "reock",
"elongation\_rectangle", "num\_holes", "hole\_area", "hole\_area\_pct",
"num\_polygons", "ew\_length", "ns\_length", "maxlength", "bearing",
"northerness", "fractaldimension", "sinuosity", "shape\_index",
"circularity\_ratio", "decimallongitude", "decimallatitude"
\end{ldescription}
\end{Arguments}
%
\begin{Details}
The function calculates various geometric metrics that describe polygon shape,
size, and orientation. For multiple features, each polygon is processed
sequentially. For parallel processing, use
\code{\LinkA{calculate\_geometric\_attributes\_parallel}{calculate.Rul.geometric.Rul.attributes.Rul.parallel}}.

\strong{Size metrics:}
\begin{itemize}

\item{} Area and perimeter provide basic size measurements
\item{} Extents (EW, NS, maxlength) measure spatial dimensions

\end{itemize}


\strong{Shape metrics:}
\begin{itemize}

\item{} Compactness measures how closely a polygon resembles a circle
\item{} Shape index and circularity provide alternative shape characterizations
\item{} Fractal dimension measures boundary complexity

\end{itemize}


\strong{Orientation metrics:}
\begin{itemize}

\item{} Bearing and northerness describe polygon orientation
\item{} Elongation quantifies directional stretching

\end{itemize}

\end{Details}
%
\begin{Value}
The input SpatVector with additional columns containing the requested
geometric attributes:
\begin{itemize}

\item{} \code{area}: Total area in square meters
\item{} \code{perimeter}: Total perimeter length in meters
\item{} \code{compactness}: Polsby-Popper compactness (1 = perfect circle)
\item{} \code{reock}: Reock compactness (area / minimum enclosing circle area)
\item{} \code{elongation\_rectangle}: Elongation ratio based on minimum bounding rectangle
\item{} \code{num\_holes}: Number of holes (interior rings) in the polygon
\item{} \code{hole\_area}: Total area of holes in square meters
\item{} \code{hole\_area\_pct}: Percentage of total area occupied by holes
\item{} \code{num\_polygons}: Number of separate polygon parts (multi-part count)
\item{} \code{ew\_length}: Average east-west extent in meters
\item{} \code{ns\_length}: Average north-south extent in meters
\item{} \code{maxlength}: Maximum distance across convex hull in meters
\item{} \code{bearing}: Geographic bearing of maximum length line in degrees
\item{} \code{northerness}: Northerness component of bearing (-1 to 1)
\item{} \code{fractaldimension}: Fractal dimension (complexity measure, typically 1-2)
\item{} \code{sinuosity}: Sinuosity (perimeter to maximum length ratio)
\item{} \code{shape\_index}: Shape index (deviation from circular shape)
\item{} \code{circularity\_ratio}: Circularity based on maximum length
\item{} \code{decimallongitude}: Centroid longitude in decimal degrees
\item{} \code{decimallatitude}: Centroid latitude in decimal degrees

\end{itemize}

\end{Value}
%
\begin{Examples}
\begin{ExampleCode}
## Not run: 
library(terra)

# Load a polygon shapefile (single or multiple features)
polygons <- vect("path/to/polygons.shp")

# Calculate all geometric attributes (processes sequentially)
polygons_with_attrs <- calculate_geometric_attributes(polygons)

# Calculate specific metrics only
polygons_subset <- calculate_geometric_attributes(polygons, 
                     metrics = c("area", "perimeter", "compactness"))

# View the results
print(polygons_with_attrs)

## End(Not run)
\end{ExampleCode}
\end{Examples}
\HeaderA{calculate\_geometric\_attributes\_parallel}{Calculate Geometric Attributes in Parallel}{calculate.Rul.geometric.Rul.attributes.Rul.parallel}
%
\begin{Description}
Calculates geometric attributes for multiple polygons. If a cluster is provided,
processes polygons in parallel; otherwise defaults to sequential processing.
\end{Description}
%
\begin{Usage}
\begin{verbatim}
calculate_geometric_attributes_parallel(v, metrics = "all", cl = NULL)
\end{verbatim}
\end{Usage}
%
\begin{Arguments}
\begin{ldescription}
\item[\code{v}] A SpatVector object representing one or more polygons.

\item[\code{metrics}] Character string or vector specifying which metrics to calculate.
See \code{\LinkA{calculate\_geometric\_attributes}{calculate.Rul.geometric.Rul.attributes}} for available options.

\item[\code{cl}] A cluster object created by \code{\LinkA{parallel::makeCluster()}{parallel::makeCluster()}}, or NULL (default).
If NULL, the function will process polygons sequentially. To enable parallel
processing, pass an explicit cluster object.
\end{ldescription}
\end{Arguments}
%
\begin{Details}
This function processes each polygon independently using the internal
single-polygon function. If a cluster is provided via the \code{cl}
parameter, it will process polygons in parallel. Otherwise, it defaults to
sequential processing (equivalent to calling \code{\LinkA{calculate\_geometric\_attributes}{calculate.Rul.geometric.Rul.attributes}}).

For a single polygon, this function simply calls
\code{\LinkA{calculate\_geometric\_attributes}{calculate.Rul.geometric.Rul.attributes}} directly.

To use parallel processing, create a cluster first:

\begin{alltt}cl <- parallel::makeCluster(parallel::detectCores() - 1)
result <- calculate_geometric_attributes_parallel(polygons, cl = cl)
parallel::stopCluster(cl)
\end{alltt}

\end{Details}
%
\begin{Value}
A SpatVector object with the same structure as input, with added
columns containing the requested geometric attributes.
\end{Value}
%
\begin{Examples}
\begin{ExampleCode}
## Not run: 
library(terra)

# Load multiple polygons
polygons <- vect("path/to/polygons.shp")

# Process sequentially (default)
result <- calculate_geometric_attributes_parallel(polygons)

# Process in parallel with explicit cluster
cl <- parallel::makeCluster(parallel::detectCores() - 1)
result <- calculate_geometric_attributes_parallel(polygons, cl = cl)
parallel::stopCluster(cl)

## End(Not run)
\end{ExampleCode}
\end{Examples}
\HeaderA{calculate\_geometric\_attributes\_single}{Calculate Geometric Attributes for a Single Polygon}{calculate.Rul.geometric.Rul.attributes.Rul.single}
%
\begin{Description}
Calculates geometric attributes for exactly one polygon feature. This is the
core computational function used by \code{\LinkA{calculate\_geometric\_attributes}{calculate.Rul.geometric.Rul.attributes}} and
\code{\LinkA{calculate\_geometric\_attributes\_parallel}{calculate.Rul.geometric.Rul.attributes.Rul.parallel}}.
\end{Description}
%
\begin{Usage}
\begin{verbatim}
calculate_geometric_attributes_single(v, metrics = "all")
\end{verbatim}
\end{Usage}
%
\begin{Arguments}
\begin{ldescription}
\item[\code{v}] A SpatVector object containing exactly one polygon feature.

\item[\code{metrics}] Character string or vector specifying which metrics to calculate.
Options are:
\begin{itemize}

\item{} \code{"all"} (default): Calculate all available metrics
\item{} A single metric name as a string: Calculate only that metric
\item{} A character vector of metric names: Calculate specified metrics

\end{itemize}

Available metric names: "area", "perimeter", "compactness", "reock",
"elongation\_rectangle", "num\_holes", "hole\_area", "hole\_area\_pct",
"num\_polygons", "ew\_length", "ns\_length", "maxlength", "bearing",
"northerness", "fractaldimension", "sinuosity", "shape\_index",
"circularity\_ratio", "decimallongitude", "decimallatitude"
\end{ldescription}
\end{Arguments}
%
\begin{Details}
This function is optimized to process a single polygon and only computes
intermediate geometric objects (convex hull, centroid, etc.) that are needed
for the requested metrics.

For processing multiple polygons, use \code{\LinkA{calculate\_geometric\_attributes}{calculate.Rul.geometric.Rul.attributes}}
(sequential) or \code{\LinkA{calculate\_geometric\_attributes\_parallel}{calculate.Rul.geometric.Rul.attributes.Rul.parallel}} (parallel).
\end{Details}
%
\begin{Value}
The input SpatVector with additional columns containing the requested
geometric attributes.
\end{Value}
%
\begin{Examples}
\begin{ExampleCode}
## Not run: 
library(terra)

# Load a single polygon
polygon <- vect("path/to/polygon.shp")[1, ]

# Calculate all attributes
result <- calculate_geometric_attributes_single(polygon)

# Calculate only specific metrics
result <- calculate_geometric_attributes_single(polygon,
  metrics = c("area", "perimeter", "compactness"))

## End(Not run)
\end{ExampleCode}
\end{Examples}
\HeaderA{calc\_elongation}{Calculate Elongation Ratio from Minimum Bounding Rectangle}{calc.Rul.elongation}
%
\begin{Description}
Calculates the elongation ratio of a polygon based on its minimum bounding
rectangle. The ratio is the major axis length divided by the minor axis length.
\end{Description}
%
\begin{Usage}
\begin{verbatim}
calc_elongation(v)
\end{verbatim}
\end{Usage}
%
\begin{Arguments}
\begin{ldescription}
\item[\code{v}] A SpatVector object representing a polygon.
\end{ldescription}
\end{Arguments}
%
\begin{Value}
A numeric value representing the elongation ratio (major/minor axis).
Higher values indicate more elongated shapes.
\end{Value}
%
\begin{Examples}
\begin{ExampleCode}
## Not run: 
library(terra)
polygon <- vect("path/to/polygon.shp")
elongation <- calc_elongation(polygon)

## End(Not run)
\end{ExampleCode}
\end{Examples}
\HeaderA{calc\_extent\_ew}{Calculate East-West Extent}{calc.Rul.extent.Rul.ew}
%
\begin{Description}
Calculates the average east-west extent of a polygon based on its bounding box.
\end{Description}
%
\begin{Usage}
\begin{verbatim}
calc_extent_ew(v)
\end{verbatim}
\end{Usage}
%
\begin{Arguments}
\begin{ldescription}
\item[\code{v}] A SpatVector object representing a polygon.
\end{ldescription}
\end{Arguments}
%
\begin{Value}
A numeric value representing the average east-west extent in meters.
\end{Value}
%
\begin{Examples}
\begin{ExampleCode}
## Not run: 
library(terra)
polygon <- vect("path/to/polygon.shp")
ew_extent <- calc_extent_ew(polygon)

## End(Not run)
\end{ExampleCode}
\end{Examples}
\HeaderA{calc\_extent\_ns}{Calculate North-South Extent}{calc.Rul.extent.Rul.ns}
%
\begin{Description}
Calculates the average north-south extent of a polygon based on its bounding box.
\end{Description}
%
\begin{Usage}
\begin{verbatim}
calc_extent_ns(v)
\end{verbatim}
\end{Usage}
%
\begin{Arguments}
\begin{ldescription}
\item[\code{v}] A SpatVector object representing a polygon.
\end{ldescription}
\end{Arguments}
%
\begin{Value}
A numeric value representing the average north-south extent in meters.
\end{Value}
%
\begin{Examples}
\begin{ExampleCode}
## Not run: 
library(terra)
polygon <- vect("path/to/polygon.shp")
ns_extent <- calc_extent_ns(polygon)

## End(Not run)
\end{ExampleCode}
\end{Examples}
\HeaderA{get\_distant\_points}{Find Most Distant Points on Convex Hull}{get.Rul.distant.Rul.points}
%
\begin{Description}
Identifies the pair of most distant points on the convex hull of a polygon,
ordered from south to north.
\end{Description}
%
\begin{Usage}
\begin{verbatim}
get_distant_points(v)
\end{verbatim}
\end{Usage}
%
\begin{Arguments}
\begin{ldescription}
\item[\code{v}] A SpatVector object representing a polygon.
\end{ldescription}
\end{Arguments}
%
\begin{Value}
A list containing:
\begin{ldescription}
\item[\code{start\_ll}] SpatVector of the southernmost point
\item[\code{end\_ll}] SpatVector of the northernmost point
\item[\code{distance}] Maximum distance in meters
\item[\code{bearing}] Geographic bearing from start to end in degrees
\end{ldescription}
\end{Value}
%
\begin{Examples}
\begin{ExampleCode}
## Not run: 
library(terra)
polygon <- vect("path/to/polygon.shp")
distant_pts <- get_distant_points(polygon)

## End(Not run)
\end{ExampleCode}
\end{Examples}
\printindex{}
\end{document}
